\lesson{3}{Mar 18 2022 Fri (16:49:12)}{Identify how make equal System of Equation}{Unit 2}
\label{les_3_unit_2:identify_how_make_equal_system_of_equation}

Let's just start with an example:

\begin{example}
  \label{exm:identify_how_to_make_equal_system_of_equation}

  \begin{equation}
    \label{cse:system_of_equation_a}
    \begin{cases}
      -4x + 3y = 4 \qquad &\textrm{\color{blue}[A1]} \\
      -5x + 6y = 14 \qquad &\textrm{\color{blue}[A2]}
    \end{cases}\qquad \textrm{System $A$}
  .\end{equation}

  \begin{equation}
    \label{cse:system_of_equation_b}
    \begin{cases}
      -4x + 3y = 4 \qquad &\textrm{\color{blue}[B1]} \\
      3x = 6 \qquad &\textrm{\color{blue}[B2]}
    \end{cases}\qquad \textrm{System $B$}
  .\end{equation}

  \begin{equation}
    \label{cse:system_of_equation_c}
    \begin{cases}
      -4x + 3y = 4 \qquad &\textrm{\color{blue}[C1]} \\
      x = 2 \qquad &\textrm{\color{blue}[C2]}
    \end{cases}\qquad \textrm{System $C$}
  .\end{equation}

  There are $3$ operations we can do to transform a system into another:

  \begin{enumerate}
    \label{enum:options_to_transform_a_system_of_equations_into_another_1}

    \item Multiply an equation by a constant
    \item Add $2$ equations together
    \item Add a multiple of one equation to another
  \end{enumerate}

  Let's see how \textbf{System A} turned into \textbf{System B}:

  We can see that the first equation in both systems stayed the same, which
  tells us that there weren't any operations done on either of them.

  However, equation $2$ clearly changed.

  Our first option is that \textbf{[A2]} was multiplied by a constant to get
  \textbf{[B2]}. But there's no number that we can multiply the first equation
  by to get the second equation. So, it's not the first option.

  Let's try our second option. If the two systems were added, it wouldn't equal
  equation \textbf{[B2]}. So, it's not the second option.

  That leaves us with our final option. Notice how in the equation
  \textbf{[B2]}, there isn't a $y$ variable, which means when we multiply the
  first equation (\textbf{[A1]}) and then add it to the second equation (
  \textbf{[A2]}), it should cancel out the $y$ variable. Now, the fist equation
  (\textbf{[A1]}) has $+3y$ and the second equation (\textbf{[A2]}) has $+6y$.
  So, to make those cancel, we should multiply the first equation
  (\textbf{[A1]}) by $-2$ then add it to the second equation (\textbf{[A2]}).

  Let's try it out:

  \begin{equation}
    \label{spt:testing_out_system_of_equation_1}
    \begin{split}
      -2(-4x + 3y) &= -2(4) \\
      8x - 6y &= -8 \\
      (-5x + 6y = 14) &+ (8x - 6y = -8) \\
      3x &= 6
    \end{split}
  .\end{equation}

  So, we were correct. To get from \textbf{System A} to \textbf{System B}, we
  leave the first equation alone, but we multiply the second equation by $-2$,
  then add the first equation to it.

  Now, how would we get from \textbf{System B} to \textbf{System C}? We leave
  the first equation alone, but you divide the second equation by $3$.
\end{example}

\newpage
